% 第 15 题 (2)：电子天平测环境温度装置示意图
% 导热汽缸（开口向上、固定桌面）内活塞 m1；轻质直杆绕支点 A，左绳挂活塞、
% 右绳挂铁块 m2 置于电子天平上。
\begin{tikzpicture}[font=\small, >=Stealth, scale=1.0]
  % 桌面
  \draw[thick, fill=gray!15] (-4.4,0) rectangle (4.4,-0.25);
  % 汽缸（开口向上，固定在桌面上）
  \draw[thick] (-3.6,0) -- (-3.6,2.6) -- (-2.2,2.6) -- (-2.2,0);
  % 活塞（m1）在缸内
  \draw[thick, fill=gray!30] (-3.5,1.15) rectangle (-2.3,1.5);
  \node[above] at (-2.9,1.5) {$m_{1}$};
  \draw[->] (-4.0,1.3) -- (-3.6,1.3);
  \node[left] at (-4.0,1.3) {$p_{0}$};
  % 支点 A 上的轻质直杆
  \draw[thick] (-3.0,3.6) -- (3.6,3.6);
  % 支点 A（三角形支撑）
  \fill (0.3,3.6) circle (1.5pt);
  \draw (0.3,3.6) -- (0.0,3.05) -- (0.6,3.05) -- (0.3,3.6);
  \node[below] at (0.3,3.05) {$A$};
  % 左细绳：杆左端 → 活塞
  \draw[thick] (-3.0,3.6) -- (-3.0,1.7);
  % 右细绳：杆右端 → 铁块
  \draw[thick] (3.6,3.6) -- (3.6,2.4);
  % 铁块 m2
  \draw[thick, fill=gray!40] (3.3,1.9) rectangle (3.9,2.4);
  \node[right] at (3.9,2.15) {$m_{2}$};
  % 电子天平
  \draw[thick] (2.6,1.15) rectangle (4.4,1.9);
  \draw[thick] (2.6,1.15) -- (2.6,0.6) -- (4.4,0.6) -- (4.4,1.15);
  \node at (3.5,1.5) {电子天平};
  % 细绳与铁块间留空（铁块放在天平上）
  \node[gray] at (3.6,2.55) {绳};
  % 标注
  \node at (-0.9,4.1) {直杆};
  \node at (-2.9,2.05) {绳};
\end{tikzpicture}
